%%%%%%%%%%%%%%%%%%%%%%%%%%%%%%%%%%%%%%%%%
% A0 Poster — Deep Food Image Classifier Project
%%%%%%%%%%%%%%%%%%%%%%%%%%%%%%%%%%%%%%%%%

\documentclass[a0,landscape]{a0poster}

% ------------------------------------------------
% Packages
% ------------------------------------------------
\usepackage{multicol}
\usepackage{graphicx}
\usepackage[svgnames]{xcolor}
\usepackage{times}
\usepackage{booktabs}
\usepackage{wrapfig}
\usepackage{amsmath}
\usepackage{caption}
\usepackage{titlesec}

% ------------------------------------------------
% Configuration
\graphicspath{{figures/}}
\columnsep=40pt
\columnseprule=3pt
\renewcommand{\columnseprulecolor}{\color{LightBlue}}
\setlength{\parskip}{2pt}
\titlespacing*{\section}{0pt}{0.4em}{0.3em}

\begin{document}
\hfill
% ------------------------------------------------
% HEADER
% ------------------------------------------------
\begin{minipage}[c]{0.19\linewidth}
    \includegraphics[width=19cm]{LTU_logo.png}
\end{minipage}
\hfill
\begin{minipage}[c]{0.60\linewidth}
\centering

{\VeryHuge\bfseries Deep Food Image Classifier with CNNs\\
and Evolutionary Hyperparameter Optimization}\\[0.3cm]
{\Large V Harsha Vardhan Yellela, MSCS Candidate \quad | \quad CJ Chung, PhD, Professor}\\
{\large College of Arts \& Science, Lawrence Technological University}
\end{minipage}
\hfill
\begin{minipage}[c]{0.19\linewidth}
    \raggedleft
    \includegraphics[width=9cm]{LTU_seal.png}
    \hspace*{5cm}
\end{minipage}
\hfill
\vspace{1cm}

\begin{multicols*}{4}

% ABSTRACT
\section*{\color{Teal}Abstract}
Food images often look similar, making classification difficult on small datasets. We evaluate a residual CNN trained on a 10‑class food dataset and explore evolutionary hyperparameter optimization methods. Evolutionary search—including ES(1+1), DEAP $(\mu{=}5,\lambda{=}10)$, CMA‑ES—and random search are compared against a Hyperband‑tuned baseline. These strategies substantially boost validation and test accuracies compared to the untuned CNN, while maintaining moderate training times.

% INTRODUCTION
\section*{\color{Teal}Introduction}
Classifying food photographs into categories such as \textit{churros}, \textit{french fries} and \textit{waffles} is challenging due to subtle visual cues, varied lighting and limited labelled samples. Deep CNNs extract hierarchical features but require careful hyperparameter tuning. Traditional manual tuning is time‑consuming; random search may miss promising regions.

\textbf{Dataset:} The simplified dataset contains 10 food classes: chicken wings, churros, french fries, hamburger, hot dog, ice cream, macaroni and cheese, pancakes, pizza, and waffles. We split 70\% of the images for training, 15\% for validation and 15\% for testing. Data augmentation applies horizontal flips, rotations up to $10^\circ$, random zoom and brightness adjustments to boost generalization on small data.

\textbf{Motivation:} Hyperparameters such as learning rate, dropout, label smoothing and augmentation strength have a large impact on CNN performance. We aim to automate this search via evolutionary algorithms.

\begin{center}
\includegraphics[width=0.95\columnwidth]{food_samples.png}
\captionof{figure}{Samples from the simplified 10‑class food dataset}
\end{center}

% SYSTEM DESIGN & WORKFLOW
\section*{\color{Teal}System Design \& Workflow}
The data flow includes image loading, preprocessing and CNN training, followed by evolutionary hyperparameter optimization. All experiments use mixed‑precision training on dual GPUs.

\begin{center}
\includegraphics[width=\columnwidth]{workflow_diagrams.png}
\captionof{figure}{Workflow: input images are resized, normalized and augmented; a ResNet‑style CNN predicts class probabilities; a confidence check triggers out‑of‑scope detection.}
\end{center}

% METHODS
\section*{\color{Teal}Methods}

\textbf{Model Architecture:}\enspace We adopt a ResNet‑style CNN with four residual blocks, expanding channel depth from 64 to 1024, followed by global average pooling and a dense head for ten classes.

\begin{center}
\includegraphics[width=\columnwidth]{cnn_architecture.png}
\captionof{figure}{ResNet‑style CNN used in all experiments}
\end{center}


\textbf{Optimization Methods:}
\vspace{-0.5em}
\begin{itemize}
    \setlength\itemsep{0.1em}
    \setlength\parskip{0.1em}
    \item \textbf{ES(1+1):} Single‑parent evolution with Gaussian mutation; mutation variance adapted via the 1/5th success rule.
    \item \textbf{DEAP $(\mu{=}5,\lambda{=}10)$:} Population‑based evolution (five parents, ten offspring) selecting the best individuals each generation.
    \item \textbf{CMA‑ES:} Covariance‑matrix adaptation to learn correlations among hyperparameters.
    \item \textbf{Random Search:} Uniform sampling across hyperparameter ranges.
    \item \textbf{Keras Tuner (Hyperband):} Bandit‑based search allocating more resources to promising configurations.
\end{itemize}
\vspace{-0.5em}
\textbf{Training:} All models are trained for up to 50 epochs with early stopping based on validation accuracy. We use cross‑entropy loss with label smoothing and Adam optimizer. Data normalization rescales pixel intensities to $[0,1]$.

\textbf{Performance Summary:}
\begin{center}
\begin{tabular}{lcccc}
\toprule
\textbf{Method} & \textbf{Val Acc (\%)} & \textbf{Test Acc (\%)} & \textbf{$\Delta$ vs Baseline} & \textbf{Macro F1} \\
\midrule
Baseline CNN & 72.33 & 72.47 & -- & 0.72 \\
ES(1+1) [Winner]      & 75.13 & 79.80 & +7.33\% & 0.80 \\
CMA‑ES       & 77.07 & 77.13 & +4.66\% & 0.77 \\
Random Search& 77.00 & 77.93 & +5.46\% & 0.78 \\
Keras Tuner  & 77.47 & 79.27 & +6.80\% & 0.79 \\
DEAP $(\mu{=}5,\lambda{=}10)$ & 77.33 & 79.33 & +6.86\% & 0.79 \\
\bottomrule
\end{tabular}
\captionof{table}{Performance summary of all methods (higher is better).}
\end{center}

\textbf{Key Observations:}  The baseline CNN underfits due to fixed hyperparameters. ES(1+1) improves accuracy by adapting learning rate, dropout and augmentation ranges over generations. CMA‑ES and DEAP both explore a larger hyperparameter space, with DEAP showing stable average improvements across populations. Random search occasionally finds good configurations but lacks directed convergence. Hyperband balances exploration and exploitation by early‑stopping poor trials and fully training promising configurations.

% CONVERGENCE BEHAVIOR
\section*{\color{Teal}Optimization Behavior}

\begin{center}
\includegraphics[width=\columnwidth]{combined_convergence.png}
\captionof{figure}{Combined convergence of ES(1+1), DEAP $(\mu{=}5,\lambda{=}10)$, CMA‑ES and Random Search. Each curve shows best‑so‑far validation accuracy.}
\end{center}

Evolutionary methods steadily improve validation accuracy. ES(1+1) converges quickly and achieves the highest performance. DEAP exhibits a gradual but stable increase. CMA‑ES explores widely with slower gains. Random Search shows noisy performance with occasional improvements.

% HYPERPARAMETER ANALYSIS
\section*{\color{Teal}Hyperparameter Analysis}

\begin{center}
\includegraphics[width=\columnwidth]{hp_comparison_grid_ansari.png}
\captionof{figure}{Hyperparameter comparison across optimization methods. Each bar shows the best value of learning rate, dropout, label smoothing, rotation and zoom found by the search.}
\end{center}

ES(1+1) and Hyperband prefer moderate learning rates and dropouts; DEAP pushes rotation and zoom augmentation further, indicating the importance of aggressive augmentation on small datasets. CMA‑ES yields balanced parameters but lower overall accuracy.

% EXPERIMENTAL RESULTS
% \section*{\color{Teal}Experimental Results}
% \begin{figure}
%     \centering
%     \includegraphics[width=0.25\linewidth]{Confusion_Matrix_(best_ES(1+1).png}
%     \caption{Enter Caption}
%     \label{fig:placeholder}
% \end{figure}
\begin{center}
\includegraphics[width=0.82\columnwidth]{Confusion_Matrix_(best_ES(1+1).png}
\captionof{figure}{Confusion matrix for the ES(1+1) optimized model on the test set, showing class-wise prediction accuracy. ES(1+1) achieves the best performance at 79.8\%.}
\end{center}

Evolutionary hyperparameter optimization consistently improves CNN performance over the untuned baseline. ES(1+1) achieves the highest test accuracy (79.8\%), followed closely by DEAP (79.3\%) and Keras Tuner (79.27\%). Population-based methods exhibit more stable validation behavior, while ES(1+1) converges faster with fewer evaluations. Detailed numerical comparisons are reported in Table 1, and convergence behavior is illustrated in Figure4.

% TRAINING CURVES & INFERENCE
\section*{\color{Teal}Training \& Inference}

To demonstrate real-world deployment behavior, we visualize inference on an unseen food image. The model outputs softmax probabilities for all classes and reports the top-5 predictions. A confidence threshold of 60% is applied to ensure prediction reliability.

\begin{center}
\includegraphics[width=\columnwidth]{inference_example.png}
\captionof{figure}{Inference output of the optimized CNN predicting churros (75.1\% confidence) with top-5 softmax probabilities and a 60\% validation threshold}
\end{center}

During inference, the trained model processes a single input image and outputs class probabilities using a softmax layer. The top-5 predictions are displayed along with confidence scores, and a confidence threshold is applied to distinguish valid food predictions from out-of-scope inputs.

% CONCLUSION
\section*{\color{Teal}Conclusion}
Evolutionary hyperparameter optimization significantly enhances CNN performance on small food classification datasets. ES(1+1) provides the highest test accuracy and converges quickly, while DEAP offers robust improvements with population diversity. Random search and CMA‑ES lag behind but still outperform the untuned baseline. Hyperband achieves competitive results but requires longer search time. The best models achieve nearly 80\% accuracy, a gain of over 7 percentage points relative to the baseline.

% FUTURE WORK
\section*{\color{Teal}Future Work}
Future directions include expanding the dataset to more classes, exploring transfer learning with pre‑trained CNN backbones, evaluating multi‑objective evolutionary strategies (accuracy and latency), and deploying models on embedded devices with quantised weights for real‑time inference.

\section*{\color{Teal}References}

\begin{enumerate}\setlength{\itemsep}{0pt}
\item I. Rechenberg, \textit{Evolution Strategy: Optimization of Technical Systems}, Frommann-Holzboog, 1973.



\item N. Hansen and A. Ostermeier, ``Completely Derandomized Self-Adaptation in Evolution Strategies,'' \textit{Evolutionary Computation}, vol. 9, no. 2, pp. 159--195, 2001. (CMA-ES)



\item L. Li, K. Jamieson, G. DeSalvo, A. Rostamizadeh, and A. Talwalkar, ``Hyperband: A Novel Bandit-Based Approach to Hyperparameter Optimization,'' in \textit{Proc. ICLR}, 2017.

\item F.-A. Fortin, F.-M. De Rainville, M.-A. Gardner, M. Parizeau, and C. Gagné, ``DEAP: Evolutionary Algorithms Made Easy,'' \textit{Journal of Machine Learning Research}, vol. 13, pp. 2171--2175, 2012.

\item K. He, X. Zhang, S. Ren, and J. Sun, ``Deep Residual Learning for Image Recognition,'' in \textit{Proc. IEEE CVPR}, 2016.

\item Kaggle Datasets, ``Food-101 Dataset,'' dansbecker/food-101.  
Available: \texttt{https://www.kaggle.com/datasets/dansbecker/food-101}
\end{enumerate}



\end{multicols*}
\end{document}
